\section{Introduction}
\label{sec:introduction}

Low Earth orbit (LEO) satellites extend Internet-of-Things (IoT)
services to remote areas, where grant-free random access (GFRA) can
reduce signaling overhead by allowing active devices to transmit
without first obtaining a grant. Since device activities are unknown
and limited pilot resources preclude mutually orthogonal pilots,
the receiver must perform joint activity detection and channel
estimation (JADCE) \cite{liu2018}, which is further complicated by
large propagation delays and Doppler shifts in LEO links.

The new radio non-terrestrial network (NR-NTN) framework in Release~17
of the 3rd Generation Partnership Project (3GPP) uses position information
from a global navigation satellite system (GNSS) receiver, satellite
ephemeris, and a broadcast common timing advance for uplink delay and
Doppler pre-compensation \cite{3gpp38300}. To support devices without
GNSS receivers or access to GNSS services, Release~19 includes a study
of GNSS-independent operation \cite{3gpp22865}. Handling differential
delays and Doppler shifts at the satellite reduces the GNSS hardware
and pre-compensation burden on IoT devices, but these offsets can
exceed one symbol duration and one subcarrier spacing, respectively.
Consequently, terrestrial orthogonal frequency division multiplexing
(OFDM)-based GFRA schemes designed for small offsets and block-fading
channels cannot be directly applied.

Waveforms for doubly selective channels offer a starting point:
orthogonal time-frequency space (OTFS) modulation exploits diversity
through a delay--Doppler representation \cite{hadani2017}, whereas
affine frequency division multiplexing (AFDM) uses a one-dimensional
discrete affine Fourier transform (DAFT) whose chirp parameters can
separate multipath components. Requiring only an inverse fast Fourier
transform and chirp multiplications, AFDM is suitable for terminals
with limited computational resources \cite{bemani2022}. The multi-frame
OTFS scheme in \cite{shen2025} collects additional observations of a
shared channel; applying this principle to AFDM requires an input--output
relationship that retains the prefix phase under long delays and
the effects of fractional delay and Doppler.

Beyond the waveform, JADCE requires a channel model suited to limited
pilot observations. Compressed sensing methods exploit sparse user
activity and channel structure, as in the joint activity and virtual
angular-domain sparsity model for massive multiple-input multiple-output
(MIMO) JADCE in \cite{ke2020}. For LEO AFDM access, a graph neural
network refines a sparse delay--Doppler prior in \cite{ding2026}.
When physical delays and Doppler shifts fall between grid points,
however, each path spreads over several coefficients. A finer grid
increases the dictionary size, while estimating continuous offsets
introduces additional unknowns coupled with the path gains.

Score-based generative models learn the channel distribution from
samples \cite{arvinte2023}. By learning the score directly on AFDM
effective channels that include fractional delay and Doppler effects,
\cite{huang2026} represents these effects without prescribing a
sparsity pattern. In multiuser GFRA, such a score-based generative
channel prior must be combined with the activity prior to recover
channels and identify active users from their superimposed signals.

In \cite{cai2026}, a turbo message-passing (MP) receiver for MIMO-OFDM
JADCE combines linear estimation, score-based channel denoising, and
an activity update without online reverse diffusion sampling. The
equivalent Gaussian observation used for denoising is an approximation
whose accuracy depends on the sensing matrix and system dimensions
\cite{cai2026}. We therefore use score pretraining to learn the channel
prior and receiver training to adapt the denoisers to the messages
encountered during multiuser inference.

Deep unfolding converts the MP iterations into a trainable network
\cite{he2020}. Since the activity probability weights the conditional
channel mean \cite{cai2026}, both estimates determine the observation
residual and the message returned to the linear estimator. An inaccurate
activity probability can thus distort interference estimation even
when the conditional channel estimate is accurate, motivating joint
training of the channel denoisers, activity classifiers, and message
parameters from received pilots.

In this paper, we investigate GNSS-independent GFRA in multiuser
MIMO-AFDM satellite IoT systems, for which we develop a multi-frame
AFDM transmission scheme and a deep-unfolded MP receiver with a
score-based generative channel prior.
The main contributions are summarized as follows.
\begin{enumerate}[label=\arabic*)]
\item We propose a multi-frame AFDM transmission scheme and derive its
input--output relationship, retaining the prefix phase for delays beyond
the useful symbol interval and including fractional delay and Doppler
in the channel operator. Under the blockwise channel model, separating
the inter-frame phase yields a multiple measurement vector (MMV) model
with one shared effective channel, so that additional pilot frames
increase observations without increasing the channel dimension.

\item We learn a score-based generative channel prior on the multiantenna
AFDM effective channel for JADCE, using samples generated with
continuous delays and Doppler shifts to account for off-grid effects.
The denoiser jointly processes each user's complete channel tensor,
exploiting dependence across DAFT input and output coordinates and
receive antennas to estimate the channel under the active hypothesis.

\item We develop a deep-unfolded MP receiver from the factorized posterior,
combining linear estimation with channel and activity denoising through
Gaussian extrinsic messages. Since the two estimates jointly determine
subsequent messages, we jointly optimize the denoisers, classifiers,
and message parameters using channel and activity supervision. Progressive
training initializes each new layer's score network from the pretrained
network and jointly fine-tunes the available layers at each depth.
\end{enumerate}

We also analyze the computational complexity of the complete receiver,
including phase and noise estimation.

% TWC-2 numerical conclusions await independent evaluation of the new receiver.

Section~\ref{sec:system} presents the system model.
Sections~\ref{sec:algorithm}--\ref{sec:overall} develop the receiver,
its training procedure, and complexity analysis. Sections~VI and VII
present numerical results and conclusions, respectively.
